This study combined hybrid sampling (SMOTEENN followed by ROS) with three
DL models (CNN, NN, MLP) and three classical ML models (DT, SVM, XGBoost)
for CVD prediction from the 2015 BRFSS survey, and evaluated them under two
protocols. When the full dataset is resampled before it is split, the
models report 94.7--99.8\% test accuracy. We showed that this performance
is an artefact of data leakage: every positive test record has an exact
copy in the training set, and the ENN step removes a fifth of the real
non-CVD records, the borderline cases, before the test set is drawn. On
6,000 real records held out before any sampling, the same models reach
only 57.4--70.8\% accuracy, and the fully grown decision tree that appeared
best falls to near chance level.

We therefore proposed a leakage-free pipeline that uses all available
features, handles class imbalance with class weights instead of synthetic
records, regularises the models and tunes the decision threshold on a
validation set. On real, unseen records it improves every model, to about
76\% accuracy, 80\% recall and an AUC of 0.84, values that five-fold
cross-validation confirms. The technical observations of this study are:

\begin{itemize}
  \item The evaluation protocol dominates the reported performance. Any
    resampling, cleaning or duplication must be applied after the test
    data have been set aside, and only to the training data; otherwise
    accuracy on imbalanced health data can be overstated by more than
    40 percentage points.
  \item On real data, hybrid sampling did not improve discrimination: it
    lowered the AUC of XGBoost from 0.823 to 0.743. Class weighting and a
    tuned decision threshold address imbalance without distorting the data.
  \item Correlation-based selection of nine features discarded useful
    information; using all 21 features raised the AUC.
  \item Model choice matters far less than data handling. Logistic
    regression, SVM, XGBoost and the three DL models all reach an AUC of
    0.83--0.84, so simple and interpretable models are a sound choice for
    this task.
  \item Compared on equal terms, the improved pipeline attains the highest
    AUC reported on the same BRFSS data (0.840), and applied unchanged to
    the clinical Cleveland data it reaches an AUC of about 0.90. Higher
    accuracies reported on clinical datasets largely reflect more
    informative clinical measurements and small, variable test sets, and
    cannot be taken as evidence of a better method without a comparable,
    leakage-free evaluation.
  \item At the real prevalence of about 10\%, a model that detects 80\% of
    CVD cases still produces three to four false alarms for every true case
    (precision about 22\%). Survey-based models are therefore suited to
    low-cost screening and risk stratification that triggers clinical
    follow-up, not to diagnosis.
\end{itemize}

Future work should validate the pipeline on external cohorts and later
BRFSS survey years, study the calibration of the predicted probabilities,
and add clinical measurements (for example blood pressure readings, lipid
profiles and ECG data), which are likely needed to move beyond the
performance ceiling of self-reported data. Explainability methods such as
SHAP values could make the models' risk factors transparent to clinicians,
and threshold selection could be aligned with the costs of missed cases
and false alarms in a specific screening programme. We also encourage
authors of CVD prediction studies to report the point at which resampling
is applied and to release their code, so that results can be verified and
compared.
